\paragraph{The retirement contract.}
\label{sec:retirement}
Fix a session $(\beta,r)$ and write $\Delta=\Delta_{i_{\beta,r}}$,
using its immutable voting snapshot. Retirement adds a local
\texttt{close\_requested} flag, closure evidence, and a first-close time
$c_v$ to the program record. These are not another phase. While draining,
an Active or Frozen program keeps that phase; only a genuine Quit or
terminal Completion changes it. The following rules extend the running
core explicitly. They do not make Force a quit proof or a terminal receipt.

\paragraph{\texttt{RequestClose}.}
The Session Manager applies this operation once per program as follows.
\begin{enumerate}[label=(\roman*),itemsep=2pt]
\item Verify complete evidence that the canonical prelude of an
authenticated finalized Espresso block closed this program's version
$\beta$. The evidence contains the triggering header, the applicable
complete control and namespace history from a validated checkpoint,
and the resulting closure and retained-prefix context. Recompute the
existing MSIG/AUTO checks; a pending MSIG, an OFF, an unfinalized header,
or a bare timestamp is insufficient. Closure does not require a selected
or Ready successor.
\item At the first such verification, atomically set
\texttt{close\_requested}, record the current time $c_v$ and the evidence,
and relay its complete validation bundle to the fixed regional membership.
This is transport of existing source proofs, not a regional vote or new
certificate. Persist a closed-version record that also prevents delayed
witnesses from creating or entering another old-$\beta$ program.
Duplicates retain the original $c_v$.
\item Close every fresh-action gate for this instance: no DATA/DRAIN
proposal, NEWVIEW, DATA/DRAIN/ANCHOR vote, new commit timer, progress
check producing BLAME, term entry, or generation entry is allowed.
The gate covers initial Enter, delayed authority/recovery callbacks,
terminal-QC rotation and buffered-message processing. No new STATUS
or other entry action can bypass it. Already dispatched messages are
not withdrawn. New-session preparation and entry use their own pair
and need not wait for this instance to retire.
\item Keep existing commit timers and their original deadlines.
Continue complete validation, comparison with retained signed material,
and evidence-first processing. A genuine current-generation quit proof
still runs Quit, broadcasts the proof, cancels its timers and preserves
commitments; it cannot reopen the entry gate. A surviving timer still
commits its block and ancestors. For each DATA thereby committed, its
Protocol~1 FINAL is delivered both to the fixed regional membership and
to the publication collector, preserving the regional dissemination rule
of \cref{def:receipt}. A partial FINAL does not invoke the complete-result
guarantee of \cref{ass:publish}.
\item Continue collecting old votes, QCs, FINAL messages and receipts,
and offer every complete result to the Publish Layer. A valid terminal
receipt still executes genuine Completion, possibly changing the phase
to Completed and canceling timers; a pending close retains its drain
deadline. Other valid receipts retain their original processing rules.
Neither receipt arrival nor a later ON reopens the closed global branch.
Keep data service and publication duties until their durable transfer.
\end{enumerate}
Messages, including valid quit evidence, are processed before a
simultaneous timer expiry or retirement event. No receipt, lock, vote,
committed block or provenance field is erased by RequestClose.

\paragraph{\texttt{FinishClose}.}
For a program that was Active or Frozen when RequestClose ran, set
\[
 d_v=c_v+4\Delta.
\]
FinishClose may mark it Terminated only at a time $t\ge d_v$ when:
\begin{enumerate}[label=(\roman*),itemsep=2pt]
\item every existing commit timer has executed or was canceled by
genuine Quit or Completion, and no queued action can pass a fresh-action
gate;
\item persistent storage owns the exit and closed-version records,
committed log, required signed evidence and validation dependencies,
receipts, anchor data and recoverable state, without releasing the last
required copy;
\item a persistent owner has accepted outstanding complete-result
offers, publication retries and data-service duties, and each locally
complete result has been offered to the Publish Layer; and
\item the dispatcher routes later messages for $(\beta,r)$ to those
owners. They may collect or relay existing signatures, validate late
receipts, serve data and continue publication, but cannot restart the
engine, execute its BFT timers or issue fresh BFT signatures for it.
\end{enumerate}
The ownership transfer may be local to the same correct replica; it
requires no additional agreement or certificate. Only after these
guards hold may the production/recovery worker and its timer machinery
be released. A late terminal receipt records archival evidence and does
not change a Terminated program back to Completed or Active.

A definitively invalid Waiting or Ready program has never entered
and has no regional timers to drain. A program already Completed when
retirement is requested has performed the original timer cancellation.
These programs may terminate after the corresponding durable exit and
responsibility transfer without the $4\Delta$ drain. A temporary change
of an unfixed candidate's eligibility alone need not terminate it.
Termination never deletes a commitment, removes DATA from $G$, or
authorizes stake withdrawal.

\begin{proposition}[Closure propagation and honest-message drain]
\label{prop:retirement-drain}
Under \cref{ass:network}, let $c=\min_{v\ {\rm correct}}c_v$ be the first
correct RequestClose time in this region, with the complete bundle
specified above. Then every correct member has closed by $c+\Delta$.
Every surviving honest commit timer expires by $c+3\Delta$, and the
regional FINAL messages generated by those timer events reach every
correct member by $c+4\Delta$. Consequently $d_v=c_v+4\Delta$ is a
conservative local drain threshold, not a deadline for finishing all
retirement duties.
\end{proposition}
\begin{proof}
The first correct relay is a complete correct-to-correct message.
Its delivery and validation are covered by the existing regional bound,
so $c\le c_v\le c+\Delta$. A correct DATA/DRAIN vote at time
$s=\tau_v(P)$ starts its only timer, with expiry $s+2\Delta$.
The close gate gives $s\le c_v\le c+\Delta$; no callback starts a
replacement timer. Thus its expiry is at most $c+3\Delta$.
Quit and Completion can cancel a timer but cannot postpone it.
The FINAL messages actually generated at surviving expiries are
delivered regionally within another $\Delta$. Finally $c_v\ge c$
implies $d_v\ge c+4\Delta$. Message-first processing covers equality.
\end{proof}

The bound starts at verified closure and its complete regional relay,
not at a control signature, the triggering Espresso timestamp, or the
start of unbounded historical catch-up. It does not add a new source
visibility assumption. Three local $\Delta$ intervals cover the last
honest timer event but not its FINAL delivery: a member can vote just
before $c+\Delta$, commit just before $c+3\Delta$, and deliver FINAL
only just before $c+4\Delta$.

Four $\Delta$ intervals do not bound subsequent receipt relays,
Byzantine message release, a storage copy, $T_{\rm pub}$, $T_E$, or
successful EXPORT acceptance. Their outstanding duties pass to the
persistent owners; no bounded storage size or transfer latency is
claimed. Elapsed grace alone therefore never authorizes FinishClose.

\begin{proposition}[Conditional regional safety during retirement]
\label{prop:retirement-safety}
Under the stated regional model and the baseline soundness premise
for DATA/DRAIN, receipts and Completion, the preceding retirement
extension preserves prefix-comparability and irrevocability of correct
regional commitments. It preserves the external prefix argument of
\cref{sec:safety-argument}. It does not preserve an unconditional
protecting-QC conclusion for every surviving timer, nor guarantee
old-session progress or complete-receipt formation.
\end{proposition}
\begin{proof}
We check the additional close gate relative to the baseline regional
soundness premise, rather than re-proving all DATA/DRAIN and Completion
rules. Consider a correct direct timer for $B$ in generation $g$,
started at $s$ and surviving to $s+2\Delta$. If the first correct
genuine Quit of $g$ were at $a\le s+\Delta$, its proof would reach
the committer by expiry and cancel this timer. The drain preserves
that processing. Therefore any such Quit satisfies $a>s+\Delta$.
Every correct entry into $g+1$ would satisfy
\[
 e_v(\beta,r,g+1,0)\ge a_v+\Delta
 \ge a+\Delta>s+2\Delta.
\]

If $s+\Delta<c$, the retirement gate suppresses no vote during
$B$'s complete-proposal propagation window. The baseline
protecting-QC and recovery-lock argument therefore applies to this
timer, with its unchanged treatment of genuine Completion.
If $s+\Delta\ge c$, its expiry is at least $c+\Delta$.
All correct members have closed by then, whereas entry into a later
generation would have to be strictly later still. The entry gate
prohibits that entry and its votes. Every later-generation QC needs
a correct signer, so missing votes during this late window cannot
be exploited through another recovery.

In one generation, complete signed histories either extend one
another or expose the original admissible equivocation proof.
A correct voter forwards its history. A correct voter on an
incompatible history either receives it before voting and detects
the proof, or forwards the incompatible history early enough to
expose the proof at the first voter by that voter's $2\Delta$
expiry. Evidence-first Quit prevents the conflicting timer
commitments. This argument uses retained evidence and continued
processing, not a new matching vote from every closed member.
Across generations, a lower-generation direct commitment cannot
belong to the late case if any correct higher-generation vote
exists. The early case's baseline protecting-QC and lock argument
then gives extension. Ancestor commitments inherit comparability;
receipt-based commitments retain their existing justification by
a correct FINAL signer, including the baseline terminal-Completion
case. No retirement step revokes a commitment.

The public interpreter still retains all of $G$ at its canonical
Force boundary and rejects further old-version DATA. Local retirement
neither selects an anchor nor authorizes an append, so it changes none
of those external transitions.
\end{proof}

The universal protecting-QC lemma from Part~I cannot simply be copied:
its proof requires matching votes from all correct members, whereas a
closed member now retains a proposal without voting. For example, with
three unit-weight members and $F=1$, let correct $H_0$ close at $c=0$
and correct $H_1$ receive closure at $\Delta$. A Byzantine leader gives
a valid proposal to $H_1$ at $0<s<\Delta$. $H_1$ votes, forwards and
later commits; $H_0$ retains it without voting. If the Byzantine member
withholds its vote and FINAL, only $H_1$'s FINAL exists. Its weight is
not greater than $F$: neither a proposal QC nor a complete receipt need
ever form. Retirement preserves this local commitment and evidence
without calling it Completion. The Publish Layer's complete-result
premise does not repair the missing signature.
